\section{Complex Structures}

Let $X$ be a topological space.

\begin{boxdefinition}[$\Delta$-complex]
    A $\Delta$-complex structure on $X$ is a set of singular $n$-simplices $\delta_{\alpha} : '\Delta^n \to X$ for varying $n = n(\alpha)$ such that
    \begin{enumerate}
        \item every $x \in X$ is the image of at least one $\delta_{\alpha}$, $\sigma_{\alpha} \vert_{\interior{\Delta^n}}$ is injective
        \item $\sigma_{\alpha}\vert_{\brac{v_0, \ldots, \widehat{v_i}, \ldots, v_n}}$ is $\sigma_{\beta} : \Delta^{k-1} \to X$ in the order-preserving identification of vertices
        \item $U \ssq X$ is open iff $\sigma_{\alpha}\inv\of{U}$ is open for all $\alpha$.
    \end{enumerate}
\end{boxdefinition}

This is the ``minimal definition'' for making the definition of homology work.

Our objective will be to show that homotopies of spaces induce homotopies of chain complexes.

\subsection{Prisms}

\begin{boxdefinition}[$n$-Prism]
    For all $n \in \N$, define the $n$-prism to be $\Delta^n \times I \ssq \R^n$.
\end{boxdefinition}

This is not a simplex, so we need to triangulate: we can view the $n$-prism as a(n uncountable) union of $(n-1)$-simplices. Denote the bottom $n$-simplex of an $n$-prism by $v_0, \ldots, v_n$ and the top $n$-simplex by $w_0, \ldots, w_n$.

\begin{boxdefinition}[Prism Operator]
    Define $P : C_n(X) \to C_{n+1}(X)$ by
    \begin{align*}
        P(\sigma) = \sum_{i} (-1)^{i} F \circ (\sigma \times \id) \vert_{\brac{v_0, \ldots, v_i, w_i, \ldots, w_n}}
    \end{align*}
\end{boxdefinition}

\begin{boxproposition}
    Given $f, g : X \to Y$, if
    \begin{align*}
        \del \circ P = g_{\#} - f_{\#} - P \circ \del
    \end{align*}
    then the maps induced by $f$ and $g$ on homology
    \begin{align*}
        f_*, g_* : H_*(X) \to H_*(Y)
    \end{align*}
    are equal.
\end{boxproposition}

\begin{boxdefinition}[Chain Homotopy]
    A sequence of maps $P_n : C_n(X) \to C_{n+1}(X)$ is called a \textbf{chain homotopy} if
    \begin{align*}
        \del \circ P + P \circ \del = g_{\#} - f_{\#}
    \end{align*}
\end{boxdefinition}

We keep in mind the following commutative diagram:
\begin{cd*}
    \cdots &[2em] {C_{n+1}(X)} &[4em] {C_n(X)} &[4em] {C_{n-1}(X)} &[2em] \cdots \\[4em]
	\cdots &[2em] {C_{n+1}(X)} &[4em] {C_n(Y)} &[4em] {C_{n-1}(X)} &[2em] \cdots
	\arrow["\partial", from=1-1, to=1-2]
	\arrow["\partial", from=1-2, to=1-3]
	\arrow["{f_{\#}}"', shift right, from=1-2, to=2-2]
	\arrow["{g_{\#}}", shift left, from=1-2, to=2-2]
	\arrow["\partial", from=1-3, to=1-4]
	\arrow["P"{description}, from=1-3, to=2-2]
	\arrow["{f_{\#}}"', shift right, from=1-3, to=2-3]
	\arrow["{g_{\#}}", shift left, from=1-3, to=2-3]
	\arrow["\partial", from=1-4, to=1-5]
	\arrow["P"{description}, from=1-4, to=2-3]
	\arrow["{g_{\#}}", shift left, from=1-4, to=2-4]
	\arrow["{f_{\#}}"', shift right, from=1-4, to=2-4]
	\arrow["\partial"', from=2-1, to=2-2]
	\arrow["\partial"', from=2-2, to=2-3]
	\arrow["\partial"', from=2-3, to=2-4]
	\arrow["\partial"', from=2-4, to=2-5]
\end{cd*}

\subsection{Exact Sequences}

Let $A \ssq X$. By $C_n(A)$, we denote the group $C_n(A; \Z)$ - the free abelian group generated by continuous functions from $\Delta^n$ to $A$.

It is, of course, a subgroup of $C_n(X)$, and since everything is abelian, all subgroups are normal, so we can construct a chain complex of quotients
\begin{align*}
    C_n\of{X, A} := \quotient{C_n(X)}{C_n(A)}
\end{align*}
\begin{boxtheorem}[Zigzag Theorem/Long Exact Sequence of Homology]\label{SP:Thm:LES_Homology}
    A short exact sequence $0\rightarrow A_{*}\rightarrow B_{*}\rightarrow C_{*} \to 0$ of chain complexes induces a long exact sequence of homology groups.
    % \begin{cd*}
    %     \cdots 
    %     \arrow[r] 
    %     & H_n(A) 
    %     \arrow[r] 
    %     & H_n(B) 
    %     \arrow[r] 
    %     & H_n(C) 
    %     \arrow[dll, 
    %          decorate,
    %          decoration={snake, amplitude=1.2pt, segment length=6pt},
    %          "\delta_n"'] \\
    %     H_{n-1}(A) 
    %     \arrow[r] 
    %     & H_{n-1}(B) 
    %     \arrow[r] 
    %     & H_{n-1}(C) 
    %     \arrow[dll, 
    %          decorate,
    %          decoration={snake, amplitude=1.2pt, segment length=6pt},
    %          "\delta_{n-1}"'] \\
    %     H_{n-2}(A) 
    %     \arrow[r] 
    %     & H_{n-2}(B) 
    %     \arrow[r] 
    %     & H_{n-2}(C) 
    %     \arrow[r] 
    %     & \cdots
    % \end{cd*}
\end{boxtheorem}
\begin{proof}
    See \cite[Week 6, Lemma 1.6]{ICLAlg4} for a proof that uses the Snake Lemma or \cite[Proposition 5.5.1]{EPFLRingsAndModules} for a more fleshed out version.
\end{proof}

Recall the following definition.

\begin{boxdefinition}
    For $A \ssq X$, we denote by $C_{\bullet}(X, A)$ the chain complex given by
    \begin{align*}
        C_n(X, A) = \quotient{C_n(X)}{C_n(A)}
    \end{align*}
    for all $n \in \N$.
\end{boxdefinition}

\begin{boxtheorem}[Long Exact Sequence of a Pair]
    For all $A \ssq X$, there is a LES
    \begin{cd*}
        \cdots \arrow[r] & H_n(A) \arrow[r] & H_n(X) \arrow[r] & H_n(X, A) \arrow[r] & H_{n-1}(A) \arrow[r] & \cdots
    \end{cd*}
\end{boxtheorem}
\begin{proof}
    Exactly \Cref{SP:Thm:LES_Homology} applied to the SES
    \begin{cd*}
        0 \arrow[r] & C_{\bullet}\of{A} \arrow[r] & C_{\bullet}\of{X} \arrow[r] & C_{\bullet}\of{X, A} \arrow[r] & 0
    \end{cd*}
\end{proof}
Say now that we have a space $X=A\cup B$ for some open $A,B$. We would like to say that $0\rightarrow C_n(A\cap B)\rightarrow C_n(A)\oplus C_n(B)\rightarrow C_n(X)\rightarrow 0$ is short exact, but this is not quite true - we can imagine a simplex $S$ which is only partially contained in each of $A,B$. However there is an easy fix for this (and it is what you would expect).

\begin{boxnotation}
    Let $X$ be a space that is covered by open sets $A$ and $B$. Define
    \begin{align*}
        C^{\set{A, B}}_n(X)
    \end{align*}
    to be the subgroup of $C_n(X)$ generated by those $\sigma : \Delta^n \to X$ such that $\pim{\sigma} \ssq A$ or $\pim{\sigma} \ssq B$.
\end{boxnotation}

The idea of the above is that we are somehow ``splitting simplices into two parts''---the part contained in $A$ and the part conained in $B$.

\begin{boxtheorem}
    Given a space $X$ and open sets $A, B \ssq X$ such that $X = A \cup B$, the following is a short exact sequence of chain complexes:
    \begin{cd*}
        0 \arrow[r] & C_{*}\of{A \cap B} \arrow[r] & C_{*}(A) \+ C_{*}(B) \arrow[r] & C^{\set{A, B}}_{*}(X) \arrow[r] & 0
    \end{cd*}
\end{boxtheorem}

Note that in the above SES, if we were to replace $C_{*}^{\set{A, B}}(X)$ by $C_{*}(X)$, the sequence would cease to be exact.

\begin{boxlemma}
    For all $n \in \N$, the $n$th homology group of $C_{*}(X)$ is the same as the $n$th homology group of $C_{*}^{\set{A, B}}(X)$.
\end{boxlemma}

More generally, we show the following.

\begin{boxlemma}
    Let $\U = \setst{U_j \ssq X}{j < J}$ be subsets of $X$ whose interiors cover $X$. Define
    \begin{align*}
        C_n^{\U}(X) := \setst{\sum_i \lambda_i \sigma_i \in C_n(X)}{\forall i,\ \exists j \st \pim{\sigma_i} \ssq U_j}
    \end{align*}
    The inclusions $i_{*} : C_{*}^{\U}(X) \inj C_{*}(X)$ form a homotopy equivalence of chain complexes.
\end{boxlemma}

By the Zigzag theorem, this gives an LES in homology:
\begin{boxdefinition}[Mayer-Vietoris Sequence]
    Given $X$ covered by the interiors of two sets $A,B$ we call the LES
    \begin{cd*}
        \cdots \arrow[r]
        &[0.1em] H_n\of{A \cap B} \arrow[r]
        & H_n(A) \+ H_n(B) \arrow[r]
        & H_n(X) \arrow[r]
        & H_{n-1}\of{A \cap B} \arrow[r]
        % & H_{n-1}(A) \+ H_{n-1}(B) \arrow[r]
        &[0.1em] \cdots
    \end{cd*}
    the \textbf{Mayer-Vietoris Sequence for $X$}.
\end{boxdefinition}

Say we have a space $X$ and $\mathcal{U}=\{U_j:j\in J\}$ is a cover of $X$ and the interiors of the $U_j$ also cover $X$. We let $C_n^{\mathcal{U}}\ssq C_n(X)$ consist of all chains $\sum_{i\in I} \lambda_i \sigma_i$ where for all $i\in I$ we have $\text{im}(\sigma_i)\ssq U_j$ for some $j$. Note that the boundary maps $\partial$ restrict to boundary maps $\partial: C_n{\mathcal{U}}\rightarrow C_{n-1}^{\mathcal{U}}(X)$. 

We now temporarily call the corresponding homology groups $H_n^{\mathcal{U}}(X)$ (``temporarily'' because we will hopefully soon show that $H_n^{\mathcal{U}}(X)$ is just $H_n(X)$ for every $\mathcal{U}$):
\begin{boxlemma}
    The inclusion map $i:C_n^{\mathcal{U}}(X)\rightarrow C_n(X)$ is a chain homotopy equivalence (i.e. the projection $p:C_n(X)\rightarrow C_n^{\mathcal{U}}(X)$ is such that $p\circ i$ and $i\circ p$ are chain homotopic to the identity), and in particular we have $H_n^{\mathcal{U}}(X)\cong H_n(X)$.
\end{boxlemma}
\begin{proof}
    We just split $\Delta^n$ into many smaller simplices, for example by repeated barycentric subdivisions... we then find a homotopy of chain complexes between $C_{*}(X)$ and $C_{*}^{\set{A, B}}(X)$.

    Let $Y \ssq \R^d$ be a convex set, and let $LC_n(Y) \ssq C_n(Y)$ denote the subgroup of affine maps $\sigma : \Delta^n \to Y$. The boundary map $\partial$ restricts to a map $\partial : LC_n(Y) \to LC_{n-1}(Y)$. For every $b \in Y$, define
    \begin{align*}
        b : LC_n(Y) \to LC_{n+1}(Y) : \brac{w_0, \dots, w_n} \mapsto \brac{b, w_0, \ldots, w_n}
    \end{align*}
    It is not difficult to show that
    \begin{align*}
        \partial \circ b = \id - b \circ \partial \iff \partial \circ b + b \circ \partial = \id
    \end{align*}

    We now let $S:LC_n(Y)\rightarrow LC_n(Y)$ be defined inductively by $S:LC_0(Y)\rightarrow LC_0(Y)$ is the identity, and $S_{\sigma}=b_{\sigma}(S\partial \sigma)$, where $B_{\sigma}$ is the barycenter of $S\partial \sigma$. Then $\partial S_{\sigma}=\partial b_{\sigma}(S\partial \sigma)=S\partial \sigma -b_{\sigma}\partial $\sorry
\end{proof}

\subsection{Application: Homology Groups of Spheres}

Let $B_n$ denote the closed unit ball in $\R^n$ centred at the origin, so that $S^{n-1} = \partial B_n \subseteq B_n$. Then, the following is a long exact sequence:
\begin{cd*}
    \cdots \arrow[r] & H_k\of{S^{n-1}} \arrow[r] & H_k\of{B_n} \arrow[r] & H_k\of{B_n, S^{n-1}} \arrow[r] & H_{k-1}\of{S^{n-1}} \arrow[r] & \cdots
\end{cd*}

We can show the following fact.

\begin{boxtheorem}
    $H_n\of{S^n; \Z} \cong \Z$ for all $n \geq 1$.
\end{boxtheorem}
\begin{proof}
    
\end{proof}

% Next up: Mayer-Vietoris

\subsection{Excision}

(Recall that relative homology is obtained by quotienting out chain complexes in $X$ by chain complexes in $A$ and then computing homology.)

\begin{boxtheorem}[Excision]
    Let $A, B \ssq X$, with the interiors covering $X$. The inclusion of $(B, A \cap B)$ into $(X, A)$ induces isomorphisms
    \begin{align*}
        H_ n\of{B, A \cap B} \xrightarrow{\sim} H_n\of{X, A}
    \end{align*}
    for all $n \in \N$. Equivalently, suppose we have a chain of inclusions $Z \ssq A \ssq X$, with $\cl{Z} \ssq \interior{A}$. Then,
    \begin{align*}
        H_n(X, A) \cong H_n\of{X \setminus Z, A \setminus Z}
    \end{align*}
    These equivalent phenomena are called \textbf{excision}.
\end{boxtheorem}
\begin{proof}
    We have a diagram
    \begin{cd*}
        \quotient{C_n(B)}{C_n(A \cap B)} \arrow[r] &
        \quotient{C_n^{\set{A, B}}(X)}{C_n(A)} \arrow[r, "\text{induces isomorphism on $H_*$}"] &[8em]
        \quotient{C_n(X)}{C_n(A)}
    \end{cd*}
\end{proof}